\section{The denominator estimate from a finite harmonic square}
\label{sec:denominator}

It remains to lower-bound \(G(z)\). This is the part where a standard
analytic presentation may invoke asymptotics for multiplicative sums.
The following direct inequality is sufficient here.

\begin{theorem}[Prime-support comparison]
\label{thm:harmonic}
For every \(s\ge1\), let \(H_s=\sum_{a=1}^s1/a\). Then
\begin{equation}
 \boxed{\qquad H_s^2\le3G(s^2).\qquad}
 \label{eq:harmonic-G}
\end{equation}
\end{theorem}

\begin{proof}
Expand the finite square:
\[
 H_s^2=\sum_{1\le a,b\le s}\frac1{ab}.
\]
Each ordered pair has the uniquely determined prime support
\(d=\rad(ab)\), which is squarefree and at most \(ab\le s^2\).
For \(d\in\sfD_{s^2}\), let
\[
 F_d(s)=
 \sum_{\substack{1\le a,b\le s\\\rad(ab)=d}}\frac1{ab}.
\]
We will show \(F_d(s)\le3h(d)\).

Write \(a=\prod_{p\mid d}p^{e_p}\),
\(b=\prod_{p\mid d}p^{f_p}\).
The exponents are nonnegative, and for each \(p\mid d\) they are not both
zero. Every exponent is at most \(s\): if \(p^e\mid a\le s\) with
\(e>0\), then \(e\le p^e\le s\); likewise for \(b\).
Uniqueness of factorization makes the map to the collection of exponent
pairs injective.

Enlarge the permitted exponent collections to the entire finite box
\[
 0\le e_p,f_p\le s,\qquad (e_p,f_p)\ne(0,0)
 \quad(p\mid d).
\]
This forgets the constraints \(a,b\le s\) and can only increase the sum
of nonnegative weights. Finite distributivity then gives
\begin{equation}
 F_d(s)\le
 \prod_{p\mid d}
 \left(
 \sum_{\substack{0\le e,f\le s\\(e,f)\ne(0,0)}}p^{-(e+f)}
 \right)
 =\prod_{p\mid d}\bigl(2A_p+A_p^2\bigr),
 \label{eq:box}
\end{equation}
where \(A_p=\sum_{e=1}^s p^{-e}\).
No infinite geometric series occurs in this step.

For the finite geometric sum, cancellation of consecutive terms gives
\[
 (p-1)A_p=1-p^{-s}\le1,\qquad
 0\le A_p\le\frac1{p-1}.
\]
If \(p\) is odd, hence \(p\ge3\), it follows that
\[
 2A_p+A_p^2
 \le\frac{2}{p-1}+\frac1{(p-1)^2}
 =\frac{2p-1}{(p-1)^2}
 \le\frac2{p-2}=h(p).
\]
The last inequality has positive denominators and reduces exactly to
\[
 2(p-1)^2-(2p-1)(p-2)=p\ge0.
\]
At \(p=2\), we instead have
\[
 2A_2+A_2^2\le3=3h(2).
\]
There is only one prime \(2\) in a squarefree support. Thus the product in
\eqref{eq:box} is at most \(h(d)\) when \(d\) is odd and at most
\(3h(d)\) when \(d\) is even. For \(d=1\), the product is empty and equals
\(1\), agreeing with the sole pair \(a=b=1\).
Hence in every case \(F_d(s)\le3h(d)\).

Finally partition the original finite pair set by \(d=\rad(ab)\):
\[
 H_s^2=\sum_{d\in\sfD_{s^2}}F_d(s)
       \le3\sum_{d\in\sfD_{s^2}}h(d)=3G(s^2).
\]
Every group with no original pair simply contributes \(0\).
\end{proof}

\begin{lemma}[Dyadic harmonic bound]
\label{lem:harmonic-dyadic}
For every \(t\ge0\),
\[
 H_{2^t}\ge1+\frac t2,\qquad
 G(2^{2t})\ge\frac{(t+2)^2}{12}.
\]
\end{lemma}

\begin{proof}
Separate the term \(1\) and the \(t\) blocks
\[
 2^{j-1}<a\le2^j,\qquad 1\le j\le t.
\]
The \(j\)-th block has \(2^{j-1}\) terms, each at least \(2^{-j}\);
its sum is at least \(1/2\). This proves the harmonic inequality by a
finite sum. Applying Theorem~\ref{thm:harmonic} with \(s=2^t\)
gives
\[
 G(2^{2t})\ge\frac{H_{2^t}^2}{3}
 \ge\frac{(1+t/2)^2}{3}
 =\frac{(t+2)^2}{12}.
\]
For \(t=0\), the calculation uses \(H_1=G(1)=1\), and remains valid.
\end{proof}
